\documentclass[a4paper]{article}
\usepackage{graphicx}
\usepackage{amsmath,amssymb}
\usepackage{minted}
\usepackage{multirow}
\usepackage{float}
\usepackage{algorithm}
\usepackage{algpseudocode}
\usepackage{todonotes}
\usepackage{forest}
\usepackage{xstring}
\usepackage{xcolor}
\usepackage[most]{tcolorbox}
\usepackage{xparse}
\usepackage{natbib}
\usetikzlibrary{graphs,graphdrawing}
\usegdlibrary{layered}
% ND Mind Mirror default input files.
\input{/home/donkarlo/Dropbox/repo/nd_language_ai_project/src/nd_language_ai/formal/computer/programming/tex/latex/code_clips/external_graphviz.tex}
\input{/home/donkarlo/Dropbox/repo/nd_language_ai_project/src/nd_language_ai/formal/computer/programming/tex/latex/part/control_sequence/kind/semtag/semtag.tex}
\input{/home/donkarlo/Dropbox/repo/nd_language_ai_project/src/nd_language_ai/formal/computer/programming/tex/latex/code_clips/seehere.tex}

\title{Theil--Sen Trend Slope Estimator}
\date{}
\author{Mohammad Rahmani}
\begin{document}
\maketitle

\section{Definition}
The Theil--Sen estimator is a robust approach to estimating the slope of a roughly linear trend. It uses the median of all valid slopes between pairs of data points rather than minimizing squared errors.

\section{Methodology}
For time points $t_i\ne t_j$, calculate the pairwise slopes
\[
 b_{ij}=\frac{y_j-y_i}{t_j-t_i},\qquad i<j,
\]
and take their median:
\[
 \widehat\beta_1^{\mathrm{TS}}
 =\operatorname{median}_{i<j} b_{ij}.
\]
One reasonable intercept convention is
\[
 \widehat\beta_0^{\mathrm{TS}}=
 \operatorname{median}_i(y_i-\widehat\beta_1^{\mathrm{TS}}t_i).
\]
Different libraries can use different intercept conventions. With irregular sampling, use actual elapsed time rather than observation numbers.

\section{Algorithm}
\begin{algorithm}[H]
\caption{Pairwise Theil--Sen slope estimation}
\begin{algorithmic}[1]
 \State $B \gets \varnothing$
 \For{$i=1,\ldots,n-1$}
  \For{$j=i+1,\ldots,n$}
   \If{$t_j\ne t_i$}
    \State Add $(y_j-y_i)/(t_j-t_i)$ to $B$
   \EndIf
  \EndFor
 \EndFor
 \State \Return $\operatorname{median}(B)$
\end{algorithmic}
\end{algorithm}

\section{E-Commerce Example}
A marketing campaign may create a single unusually high revenue day. Such a point can strongly influence a least-squares slope. The median pairwise slope is less affected by a small number of outliers and can provide a more stable description of underlying sales growth.

\section{Advantages}
\begin{itemize}
 \item More resistant to isolated outliers than ordinary least squares.
 \item Direct interpretation in revenue per day or orders per week.
 \item Does not require Gaussian errors to define a point estimate.
\end{itemize}

\section{Limitations}
\begin{itemize}
 \item A naive implementation computes $O(n^2)$ slopes and may use $O(n^2)$ memory.
 \item A single global slope misses reversals and structural breaks.
 \item Strong seasonality and time-series dependence still require care in interpretation and uncertainty estimates.
\end{itemize}

\section{Practical Guidance}
Use this estimator for magnitude of a trend and a suitable statistical test for evidence of its presence. Compare results with a seasonally adjusted linear model when sales show a strong weekly pattern.

\bibliographystyle{plainnat}
\bibliography{/home/donkarlo/Dropbox/repo/nd_shared_research/bibliography/bibliography.bib}
\end{document}
